\documentclass[10pt,a4paper]{amsart}
\usepackage[margin=19mm]{geometry}
\usepackage{amsmath,amssymb,mathtools}
\usepackage[scr=rsfs,cal=euler]{mathalfa}
\usepackage{xfrac}
\usepackage[unicode,hidelinks,bookmarks=false]{hyperref}
\newcommand{\ie}{i.e.\ }
\newcommand{\cf}{cf.\ }
\newcommand{\resp}{respectively\ }
\newcommand{\Subsection}{\subsection}
\providecommand{\nobreakdash}{\nobreak\hbox{-}}
\setlength{\emergencystretch}{2em}
\multlinegap0pt
\usepackage{bm}
\usepackage{bbm}
\usepackage{dsfont}
\usepackage{xspace}
\usepackage{cancel}
\usepackage{mathtools}
\usepackage[shortlabels]{enumitem}
\usepackage{amsthm}
\makeatletter
\@ifundefined{theorem}{\newtheorem{theorem}{Theorem}}{}
\@ifundefined{proposition}{\newtheorem{proposition}[theorem]{Proposition}}{}
\makeatother
\DeclareMathOperator*{\argmin}{arg\,\min}
\DeclareMathOperator{\supp}{Supp}
\newcommand{\R}{\mathbb{R}}
\newcommand{\di}{\mathrm{d}}
\title[Quantitative Stability in Discrete Optimal Transport]{Quantitative Stability in Discrete Optimal Transport}
\author{William Ford}
\date{Source: arXiv:2510.17407v1 (CC BY 4.0)}
\begin{document}
\maketitle

\noindent\emph{Source and licence.} Quantitative Stability in Discrete Optimal Transport. Version arXiv:2510.17407v1 (CC BY 4.0), distributed under the Creative Commons Attribution 4.0 International licence, as recorded in the arXiv licence metadata for this exact version and shown on the abstract page https://arxiv.org/abs/2510.17407v1. Full rights notice: licenses/arxiv-2510.17407-CC-BY-4.0.txt.\par\smallskip

\noindent\textit{Editorial scope.} Counted unit: the Theorem whose source label is \texttt{thrm: structure of discrete plans new notation} in the article (source0.tex lines 1620--1635, source label \texttt{thrm: structure of discrete plans new notation}), delivered together with the article's own complete original proof of it (source0.tex lines 1636--1670, one \texttt{proof} environment of 35 lines). No mathematics is written, repaired or re-derived: statement and proof are the article's own text line for line. Required context retained uncounted: eq: discrete transport problem X and Y (lines 961--964); eq: distinct points (lines 975--978); prop: uniqueness iff strict cm (lines 989--996). Every definition, display and label of those lines is retained unchanged. Omitted from this package and not redistributed: 1--960, 965--974, 979--988, 997--1619, 1671--2508. Nothing inside a retained excerpt is omitted or altered: every retained body is delivered exactly as the article's lines are written, so no omission receipt of any kind accompanies this package. Citations: the retained excerpts carry no citation command at all, so no bibliography entry of the article is transcribed and no reference is redistributed. Reference adaptation: none was needed and none was made; every reference inside the delivered unit resolves inside it, so no unresolved reference or citation is produced. Printed slips kept literal: no printed word, symbol, hypothesis or number of the counted statement or of its proof was altered. Layout and class: the article's own class is \texttt{book} with packages including inputenc, geometry, adjustbox, changepage, xargs, xcolor, todonotes, titlesec, setspace, graphics, tocbibind, graphicx, amsmath, amsfonts; the wrapper is the lane's pinned \texttt{amsart} editorial wrapper at 10pt on A4, which restates the theorem-like environments the retained text needs and adds the article's own macro definitions that the retained text needs (\texttt{\textbackslash R}, \texttt{\textbackslash argmin}, \texttt{\textbackslash di}, \texttt{\textbackslash supp}), with extra wrapper packages (bm, bbm, dsfont, xspace, cancel, mathtools, enumitem). The delivered numbering is the unit's own. Assets: no retained span contains a figure environment, a graphics-inclusion command, a PDF-page inclusion, a table or a TikZ picture; no image file is copied into this package, the package declares no asset dependency and no image file is redistributed with it. Licence: version arXiv:2510.17407v1 of this article is distributed under the Creative Commons Attribution 4.0 International licence, as recorded in the arXiv licence metadata for this exact version and shown on the abstract page https://arxiv.org/abs/2510.17407v1. A full rights notice is delivered at licenses/arxiv-2510.17407-CC-BY-4.0.txt. Characters: every retained excerpt is pure ASCII, so the delivered engine, XeLaTeX with its default Latin Modern text font, reports no missing character.
\begin{equation}
\label{eq: discrete transport problem X and Y}
    \Gamma_c(X, Y) : = \argmin_{\gamma \in \Pi(\rho_X, \mu_Y)} \int_{\R^d \times \R^d} c(x, y) \di \gamma(x, y),
\end{equation}

\begin{equation}
\label{eq: distinct points}
  x_i \neq x_j \quad\text{ and }\quad y_i \neq y_j \quad \text{ for all } i \neq j.  
\end{equation}

\begin{proposition}
\label{prop: uniqueness iff strict cm}
    Let $X \in (\R^d)^M$ and $Y \in (\R^d)^N$ and consider the transport problem \eqref{eq: discrete transport problem X and Y} for some $c: \R^d \times \R^d \to \R$. Then
    \begin{enumerate}[label=(\roman*)]
        \item \label{enum: weak inequality} $\gamma \in \Gamma(X, Y)$ if and only if $\Delta_P \geq 0$ for all distinct (in the sense of \eqref{eq: distinct points}) $P$ in $\supp \gamma$.
        \item \label{enum: strict gives uniqueness} Problem \eqref{eq: discrete transport problem X and Y} has uniqueness of solutions if and only if there exists $\gamma \in \Pi(\rho_X, \mu_Y)$, such that $\Delta_P > 0$ for all distinct (in the sense of \eqref{eq: distinct points}) $P \subset \supp \gamma$ of length $\geq 2$ (and this $\gamma$ is the unique solution).
    \end{enumerate}
\end{proposition}

\begin{theorem}[Structure of discrete plans along trajectories preserving uniqueness]
\label{thrm: structure of discrete plans new notation}
Let $c: \R^d \times \R^d \to \R_+$ be continuous. Let $X: [0, 1] \to (\R^d)^M$ and $Y : [0, 1] \to (\R^d)^N$ be continuous curves such that
\begin{enumerate}[label = (\roman*)]
    \item \label{enum: distinct points (0, 1)} For each $t \in (0, 1)$ and any distinct indices $i \neq j$, $x_i(t) \neq x_j(t)$ and $y_i(t) \neq y_j(t)$.
    \item \label{enum: unique solutions for each pb} For each $t \in [0, 1]$ solutions to the $c$-cost transport between $\rho_{X(t)}$ and $\mu_{Y(t)}$ are unique.
\end{enumerate}
Then there exists $\hat{\gamma} \in \Pi(\alpha, \beta)$ such that
\begin{equation*}
    \Gamma_c(X(t), Y(t)) = \left\{ \sum_{i = 1}^{M}\sum_{j=1}^N \hat\gamma_{ij} \delta_{x_i(t), y_j(t)} \right\} \text{ for all } t \in [0, 1],
\end{equation*}
so that the ``same" plan is the unique optimal plan for all $t \in [0, 1]$. In other words, given the optimal plan for some $s \in (0, 1)$, the unique optimal plan for each other $t \in [0, 1]$ is a glue-composition (note there can be multiple glue compositions here, but only one can be optimal since we assume uniqueness) of $ \pi_{ts}^\rho \circ \gamma(s) \circ \pi_{st}^\mu$ where
\begin{equation*}
    \pi_{ts}^\rho := \sum_{i = 1}^M \alpha_i \delta_{x_i(t), x_i(s)} \in \Pi(\rho_{X(t)}, \rho_{X(s)})\; \text{ and } \; \pi_{st}^\mu :=\sum_{i = 1}^M \beta_i \delta_{y_i(s), y_i(t)} \in \Pi(\mu_{Y(s)}, \mu_{Y(t)}).
\end{equation*}
\end{theorem}

\begin{proof}
Fix some $s \in (0, 1)$ and let $\gamma(s)$ be the unique optimal plan between $\rho_{X(s)}$ and $\mu_{Y(s)}$ for cost $c$, written in the form
\begin{equation}
\label{eq: gamma(s) definion}
    \gamma(s) = \sum_{i = 1}^{M}\sum_{j=1}^N \hat\gamma_{ij}(s) \delta_{x_i(s), y_j(s)}. \quad \text{ (Here some $\hat\gamma_{ij}(s)$ may be zero.)}
\end{equation}
Here a priori, the $\hat\gamma_{ij}(s)$ \textit{depend on the choice of $s$}. We want to show that they are independent of $s$, so that the same $\gamma(t)$ is optimal for all  $t \in [0, 1]$, rather than just $t = s$. The ways to choose a finite number of ways to choose distinct $\hat P \subset \hat\gamma(s)$ and distinct $P \subset \gamma(t)$ are in bijection, and there are a finite number of possibilities $r$. We denote these by $\{P_k(t)\}_{k=1}^r$, where each is of the form
\begin{equation*}
    P_k(t) = \{(x_{i_l}(t), y_{j_l}(t))\}_{l=1}^{L_k} \quad \text{ with }\quad x_{i_{l}}(t) \neq x_{i_{l'}}(t), y_{j_l}(t) \neq y_{j_{l'}}(t)\; \forall l \neq l', t \in (0, 1).
\end{equation*}
By Proposition \ref{prop: uniqueness iff strict cm}, $\Delta_{P_k(s)} >0$, and since $c$, $X(t)$ and $Y(t)$ are continuous, so is $t \mapsto \Delta_{P_k(t)}$. As there are a finite number of distinct support selections,
\begin{equation*}
    \underline{\Delta}(t) : = \inf_{k = 1,..., r} \Delta_{P_k(t)}
\end{equation*}
is continuous and $\underline{\Delta}(s)>0$. Set
\begin{equation*}
    t^+ = \inf \left\{ t \geq s : \underline{\Delta}(t) = 0 \right\} \quad \text{ and } t^- = \sup\left\{t \leq s : \underline{\Delta}(t) = 0\right\}.
\end{equation*}
Assume $t^- \in (0, s)$, then there exists $P \subset \supp \gamma(t^-)$ such that $\Delta_P = 0$, and by continuity and definition of $t^-$ as a supremum, $\Delta_{P'} \geq 0$ for all other $P' \subset \supp \gamma(t^-)$. Proposition \ref{prop: uniqueness iff strict cm} implies both that $\gamma(t^-)$ is optimal and that it is not the only optimiser, contradicting the non-uniqueness assumption. Hence, it must be that $\underline{\Delta}(t) > 0$ for all $t \in (0,s)$. Arguing symmetrically for $t^+$, we deduce that $\underline{\Delta}(t) > 0$ for all $t \in (0,1)$, and hence is the unique optimiser. This extends to $t \in [0, 1]$ by qualitative stability.

To explicit the glueing structure, for $s \in (0, 1)$ and $t \in [0, 1]$ consider the measure
\begin{equation}
\label{eq: G(s, t) definition glueing measure}
    \eta(s,t) := \sum_{i=1}^M \sum_{j = 1}^N \hat\gamma_{ij} \delta_{x_i(t), x_i(s), y_j(s), y_j(t)} \in \Pi(\rho_{X(t)}, \rho_{X(s)}, \mu_{Y(s)}, \mu_{Y(t)}).
\end{equation}
It follows that
\begin{equation*}
    p_{12\#}\eta(s, t) = \pi^\rho_{ts};\quad p_{23\#} \eta(s, t) = \gamma(s); \quad \text{ and } p_{34\#} \eta(s, t) = \pi^\mu_{st}
\end{equation*}
and
\begin{equation*}
    p_{14\#} \eta(s, t) = \gamma(t) \quad\forall t \in [0, 1].
\end{equation*}
Thus, we have shown that the optimal plan over $\Pi(\rho_t, \mu_t)$ is given by glue composition, supported by the measure $\eta(s, t)$.
\end{proof}

\end{document}
